\section{Hierarki Chomsky Tipe 3 sampai Tipe 0}
\label{sec:ch04-hierarki}

Noam Chomsky mengklasifikasikan tata bahasa berdasarkan bentuk aturan
produksinya. Setiap tipe bersesuaian dengan kelas bahasa dan mesin
pengenal tertentu. Semakin longgar bentuk aturan, semakin besar daya
ekspresif bahasa yang dihasilkan. Tabel dan urutan pembahasan pada
bagian ini mengikuti Guntoro Bab 6, ``Hirarki Chomsky''.

\begin{figure}[H]
\centering
\begin{tikzpicture}
  \draw[thick, rounded corners=4pt, fill=ubbblue!5]
        (0,0) rectangle (13.2,7.2);
  \node[anchor=north west, font=\small\bfseries] at (0.2,7.0)
        {Tipe 0: rekursif terbilang (mesin Turing)};
  \draw[thick, rounded corners=4pt, fill=ubbblue!10]
        (0.7,0.7) rectangle (12.5,6.0);
  \node[anchor=north west, font=\small\bfseries] at (0.9,5.8)
        {Tipe 1: sensitif konteks (otomat terbatas linear)};
  \draw[thick, rounded corners=4pt, fill=ubbblue!15]
        (1.4,1.4) rectangle (11.8,4.8);
  \node[anchor=north west, font=\small\bfseries] at (1.6,4.6)
        {Tipe 2: bebas konteks (otomat pushdown)};
  \draw[thick, rounded corners=4pt, fill=ubbblue!25]
        (2.1,2.1) rectangle (11.1,3.6);
  \node[anchor=north west, font=\small\bfseries] at (2.3,3.4)
        {Tipe 3: regular (otomat hingga)};
\end{tikzpicture}
\caption{Hierarki Chomsky sebagai kelas bahasa yang bersarang}
\label{fig:ch04-hierarki}
\end{figure}

\subsection{Tipe 3: Tata Bahasa Regular}

Tata bahasa regular memiliki aturan berbentuk
\[
A \to aB, \quad A \to a, \quad A \to \varepsilon,
\]
dengan \(A, B \in V\) dan \(a \in T\). Bentuk ini disebut
\emph{right-linear}. Bentuk \emph{left-linear} \(A \to Ba\) juga
menghasilkan kelas bahasa yang sama. Tata bahasa regular ekivalen
dengan ekspresi reguler dan otomata hingga.

Contoh tata bahasa regular untuk identifier:
\[
\begin{aligned}
I &\to aI \mid bI \mid \ldots \mid zI \mid A \mid \ldots \mid Z
      \mid \_I,\\
A &\to aA \mid \ldots \mid zA \mid 0A \mid \ldots \mid 9A
      \mid \varepsilon.
\end{aligned}
\]
Tata bahasa ini membangkitkan identifier yang dimulai dengan huruf
atau garis bawah, diikuti huruf, angka, atau garis bawah.

\subsection{Tipe 2: Tata Bahasa Bebas Konteks}

Tata bahasa bebas konteks (CFG) memiliki aturan \(A \to \alpha\),
dengan \(A \in V\) dan \(\alpha \in (V \cup T)^{*}\). CFG
membangkitkan bahasa bebas konteks, yang dikenali oleh otomata
pushdown.

Contoh CFG untuk ekspresi aritmetika:
\[
E \to E + T \mid T, \quad
T \to T * F \mid F, \quad
F \to (E) \mid \code{id}.
\]
CFG juga dapat membangkitkan bahasa dengan struktur bersarang,
misalnya \(S \to aSb \mid \varepsilon\) untuk
\(\{a^n b^n \mid n \ge 0\}\).

Kelas CFG tertutup terhadap gabungan, konkatenasi, dan bintang
Kleene, tetapi tidak tertutup terhadap irisan dan komplemen. Sebagai
contoh, irisan \(\{a^n b^n c^m\}\) dan \(\{a^m b^n c^n\}\) adalah
\(\{a^n b^n c^n\}\), yang bukan bebas konteks.

\subsection{Tipe 1: Tata Bahasa Sensitif Konteks}

Tata bahasa sensitif konteks memiliki aturan
\[
\alpha A \beta \to \alpha \gamma \beta,
\]
dengan \(\gamma \neq \varepsilon\). Aturan hanya boleh mengganti
nonterminal \(A\) dalam konteks \(\alpha\) dan \(\beta\), dan panjang
string tidak boleh berkurang. Kelas ini dikenali oleh otomata
terbatas linear (\eng{linear bounded automaton}).

Contoh bahasa sensitif konteks adalah
\(\{a^n b^n c^n \mid n \ge 1\}\). Bahasa ini memerlukan pencocokan
tiga kelompok dengan jumlah sama, yang tidak dapat dilakukan oleh
satu stack.

\subsection{Tipe 0: Tata Bahasa Tanpa Batas}

Tata bahasa tanpa batas (\eng{unrestricted grammar}) memiliki aturan
\(\alpha \to \beta\) tanpa batasan bentuk, asalkan \(\alpha\) memuat
paling sedikit satu nonterminal. Kelas ini ekivalen dengan mesin
Turing dan mencakup bahasa yang dapat dikenali secara rekursif
terbilang.

Sebagian bahasa pada kelas ini tidak dapat diputuskan. Contoh
klasiknya adalah masalah penghentian (\eng{halting problem}): tidak
ada algoritma yang dapat menentukan apakah program sembarang berhenti
untuk semua masukan.

\begin{table}[H]
\centering
\caption{Ringkasan hierarki Chomsky}
\label{tab:ch04-ringkasan}
\begin{tabularx}{\textwidth}{@{}l l X l@{}}
\toprule
\textbf{Tipe} & \textbf{Kelas bahasa} & \textbf{Bentuk aturan} &
\textbf{Automata} \\
\midrule
3 & Regular & \(A \to aB\), \(A \to a\) & DFA/NFA \\
2 & Bebas konteks & \(A \to \alpha\) & PDA \\
1 & Sensitif konteks & \(\alpha A\beta \to \alpha\gamma\beta\) &
Otomata terbatas linear \\
0 & Rekursif terbilang & \(\alpha \to \beta\) & Mesin Turing \\
\bottomrule
\end{tabularx}
\end{table}

\subsection{Menyusun Tabel Pemetaan}

Tabel \tabref{tab:ch04-ringkasan} adalah artefak Sub-CPMK 1.4. Untuk
melengkapinya, tuliskan pada setiap baris: tipe, kelas bahasa, bentuk
aturan, jenis automata, dan satu contoh bahasa. Contoh yang dapat
dipakai adalah \(a^{*}b^{*}\) untuk Tipe 3, \(\{a^n b^n\}\) untuk
Tipe 2, \(\{a^n b^n c^n\}\) untuk Tipe 1, dan bahasa yang dikenali
mesin Turing untuk Tipe 0. Periksa bahwa setiap contoh benar-benar
termasuk kelas yang dipetakan, bukan hanya mirip secara intuisi.
